% Generated draft: requires mathematical and source review.
\documentclass[11pt]{article}
\usepackage[margin=25mm]{geometry}
\usepackage{fontspec}
\setmainfont{Latin Modern Roman}
\usepackage{amsmath,amssymb,mathtools,unicode-math}
\setmathfont{Latin Modern Math}
\usepackage{enumitem,needspace,xurl,hyperref}
\setlength{\parindent}{0pt}
\setlength{\parskip}{5pt}
\setlength{\emergencystretch}{3em}
\newcommand{\N}{\mathbb{N}}
\newcommand{\Z}{\mathbb{Z}}
\newcommand{\Q}{\mathbb{Q}}
\newcommand{\R}{\mathbb{R}}
\newcommand{\C}{\mathbb{C}}
\newcommand{\F}{\mathbb{F}}
\newcommand{\A}{\mathbb{A}}
\newcommand{\PP}{\mathbb{P}}
\newcommand{\E}{\mathbb{E}}
\newcommand{\eps}{\varepsilon}
\newcommand{\dd}{\,\mathrm{d}}
\begin{document}
\section*{UnsolvedMath ID: 30004221}
\textbf{Source problem code:} catalogue entry
\subsection{Definitions and mathematical statement}
Let \(V=\C^n\). Relative to the representation-theoretic setting intended by the source, let \(\lambda\) be a dominant weight and let \(\mu\) be a weight of \(\Lambda^k V\) such that \(\lambda+\mu\) is dominant. Let \(\Phi^+\) denote the intended set of positive roots, let \(\rho\) denote the intended Weyl vector, and write \(\langle\cdot,\cdot\rangle\) for the pairing used in the source. Define
\[
S_{\lambda,\mu}=\left\{\alpha\in\Phi^+:
\langle\lambda+\rho+\mu,\alpha\rangle
=\langle\lambda+\rho,\alpha\rangle-1\right\}.
\]
The proposed identity, for every admissible choice of these data, is
\[
\kappa_{\lambda,\mu}
\stackrel{?}{=}
\prod_{\alpha\in S_{\lambda,\mu}}
\frac{\langle\lambda+\rho,\alpha\rangle}
{\langle\lambda+\rho+\mu,\alpha\rangle}.
\]
Here \(\kappa_{\lambda,\mu}\) is described in the supplied record as the resulting one-dimensional local intersection form, or local-intersection number. The empty product is interpreted as \(1\).

This is a qualified transcription, not a fully specified numerical conjecture. The record does not identify the acting group or root datum, specify the root-pairing and dominance conventions, define \(\rho\) explicitly, or give the construction of the local intersection form. In particular, a one-dimensional bilinear form does not determine a numerical coefficient without a specified normalization or basis convention. The displayed fractions also presuppose nonzero denominators; the supplied record does not explicitly establish this. The admissible ranges of \(n\) and \(k\) are not stated. These omissions must be resolved before the identity can be interpreted unambiguously.
\subsection{Short English statement}
The proposed clasp formula expresses a local intersection number associated with \(\lambda\) and a weight \(\mu\) of \(\Lambda^k\C^n\) as a product over positive roots whose pairing decreases by exactly \(1\) after adding \(\mu\). The intended representation-theoretic conventions and the normalization of the intersection number require clarification. The supplied open-status assessment concerns a different conjecture and does not establish the status of this formula.
\subsection{Sources}
\begin{itemize}
\item Local registry: \texttt{data/unsolvedmath/problems.json}, permanent UnsolvedMath ID \(30004221\), title ``Clasp Formula for Local Intersection Numbers,'' exact source problem code \texttt{catalogue entry}. The supplied proposed year is 2019; no proposer is identified.
\item Supplied source citation: \emph{Geometric, Algebraic, and Topological Combinatorics} (2020), \emph{Oberwolfach Reports}, pp.\ 2395--2472, DOI \texttt{10.4171/owr/2019/39}. \url{https://doi.org/10.4171/owr/2019/39}. This report is also cited in the background as documenting workshop open-problem discussions.
\item Supplied literature source: Adler, De Loera, Klee, and Zhang, ``cocircuit diameters.'' \url{https://doi.org/10.37236/9653}.
\item Supplied literature source: ``Oberwolfach Report 39, oriented matroids.'' \url{https://doi.org/10.4171/owr/2019/39}.
\item The literature assessment is dated exactly \texttt{2026-08-22}. It is supplied dated curation, not independently verified here. It labels the problem open but discusses a quadratic bound, counterexamples to stronger variants, and an \(n-r+2\) conjecture. That discussion is not matched to the local intersection formula in this record.
\item The registry labels the statement \texttt{corrected\_verified} and reports that the official Oberwolfach report was checked directly. This is a supplied verification claim, not an independent source check performed for this draft. The supplied citation is retained, but its support for the exact formula and its conventions requires human verification.
\end{itemize}
\end{document}
